凸优化给出多项式时间可解的问题类，但现实中有大量问题是 NP-hard 的：整数规划、旅行商、图着色、最大割。面对这些问题，精确求解在计算上不可行，只能寻找近似解。计算复杂度理论告诉我们哪些问题本质困难，近似算法则告诉我们困难问题能近似到什么程度。

本单元回答三个问题：

\begin{enumerate}
\tightlist
\item
  P、NP 与 NP-hard 怎样给问题分类：多项式时间可解、多项式时间可验证、以及``至少和所有 NP 问题一样难''的含义。
\item
  凸松弛把 NP-hard 问题变成可解问题：整数规划的 LP 松弛、秩约束的核范数松弛、最大割的 SDP 松弛各自保留了多少结构。
\item
  近似算法给出最坏情况保证：近似比的定义、几个经典结果（最大割 0.878、旅行商 1.5）、以及近似 hardness 的边界。
\end{enumerate}

这三个问题共用一条主线：计算可行性是比最优性更基本的约束。一个算法在最坏情况下需要指数时间，即使它总是给出最优解，在实践中也可能不如一个总是给出 90\% 最优解的多项式算法。

\subsection{一个图例区分三种难度}

以三顶点的完全图 \(K_3\) 为例。旅行商问题要求找一条访问每点恰好一次并回到起点的最短回路。3 个顶点的所有回路只有 \(3!/3=2\) 条，穷举即可。当顶点增加到 100，回路数量为 \(99!/2\approx4.6\times10^{155}\)，超出了任何可计算穷举。这就是输入尺寸 \(n\) 增长时问题难度的剧变。对于 LP 可解的 assignment 松弛，\(n=100\) 时最优值 342，而 ILP 最优值 379，间隙 37。松弛给出了 342 的下界，舍入到 379 后，我们证明了 379/342=1.108，即至多比真实最优大 11\%——计算可行性、松弛质量和近似比在同一个小例中紧密耦合。

\subsection{怎么读这一章}

先区分``数学上有最优解''和``能在可接受时间内找到它''。第一篇建立 P、NP 与 NP-hard 的语言，第二篇再看遇到难问题时怎样用凸松弛取得可算的下界和近似解。重点不是背复杂度类别，而是学会说明算法究竟牺牲了什么。非凸算法能提供的保证可回看 \hyperref[ch:078]{``非凸与随机优化：知道算法真正保证了什么''}。

\subsection{本章篇目}

以下篇目按概念依赖排列。

\begin{itemize}
\tightlist
\item \hyperref[sec:09-Convex-Optimization/09-Complexity-and-Approximation/01]{``P、NP 与 NP-hard：给问题按计算难度分类''}；
\item \hyperref[sec:09-Convex-Optimization/09-Complexity-and-Approximation/02]{``凸松弛把 NP-hard 问题变成可解问题''}。
\end{itemize}
